%-------------------------
% Resume in LaTeX
% Based on: https://github.com/jakegut/resume
% License: MIT
%-------------------------

\documentclass[letterpaper,11pt]{article}

\usepackage{latexsym}
\usepackage[empty]{fullpage}
\usepackage{titlesec}
\usepackage{marvosym}
\usepackage[usenames,dvipsnames]{color}
\usepackage{verbatim}
\usepackage{enumitem}
\usepackage[hidelinks]{hyperref}
\usepackage{fancyhdr}
\usepackage[english]{babel}
\usepackage{tabularx}
\input{glyphtounicode}

\pagestyle{fancy}
\fancyhf{}
\fancyfoot{}
\renewcommand{\headrulewidth}{0pt}
\renewcommand{\footrulewidth}{0pt}

% Adjust margins
\addtolength{\oddsidemargin}{-0.5in}
\addtolength{\evensidemargin}{-0.5in}
\addtolength{\textwidth}{1in}
\addtolength{\topmargin}{-.5in}
\addtolength{\textheight}{1.0in}

\urlstyle{same}

\raggedbottom
\raggedright
\setlength{\tabcolsep}{0in}

% Section formatting
\titleformat{\section}{
  \vspace{-5pt}\scshape\raggedright\large
}{}{0em}{}[\color{black}\titlerule \vspace{-5pt}]

% ATS-friendly: ensure PDF is machine-readable
\pdfgentounicode=1

%-------------------------------------------
% Custom commands
\newcommand{\resumeItem}[1]{
  \item\small{#1 \vspace{-2pt}}
}

\newcommand{\resumeSubheading}[4]{
  \vspace{-2pt}\item
    \begin{tabular*}{0.97\textwidth}[t]{l@{\extracolsep{\fill}}r}
      \textbf{#1} & #2 \\
      \small#3 & \small #4 \\
    \end{tabular*}\vspace{-7pt}
}

\newcommand{\resumeProjectHeading}[2]{
    \item
    \begin{tabular*}{0.97\textwidth}{l@{\extracolsep{\fill}}r}
      \small#1 & #2 \\
    \end{tabular*}\vspace{-7pt}
}

\newcommand{\resumeSubItem}[1]{\resumeItem{#1}\vspace{-4pt}}

\renewcommand\labelitemii{$\vcenter{\hbox{\tiny$\bullet$}}$}

\newcommand{\resumeSubHeadingListStart}{\begin{itemize}[leftmargin=0.15in, label={}]}
\newcommand{\resumeSubHeadingListEnd}{\end{itemize}}
\newcommand{\resumeItemListStart}{\begin{itemize}}
\newcommand{\resumeItemListEnd}{\end{itemize}\vspace{-5pt}}

%===========================================================================
%                             DOCUMENT START
%===========================================================================
\begin{document}

%-----------HEADING-----------
\begin{center}
    \textbf{\Huge \scshape Oleksii Krasnoshtanov} \\ \vspace{1pt}
    \small
    \href{mailto:oleksiikrasnoshtanov@gmail.com}{\underline{oleksiikrasnoshtanov@gmail.com}} $|$
    +31 6 29746862 $|$
    Tilburg, Netherlands $|$
    \href{https://linkedin.com/in/oleksii-krasnoshtanov}{\underline{LinkedIn}} $|$
    \href{https://github.com/alex-krasnoshtanov}{\underline{GitHub}}
\end{center}

%-----------EDUCATION-----------
\section{Education}
  \resumeSubHeadingListStart
    \resumeSubheading
      {Breda University of Applied Sciences}{Breda, Netherlands}
      {Bachelor of Science in Applied Data Science \& Artificial Intelligence; 9.1/10 (4.0 GPA)}{Sep 2024 -- Present}
      \resumeItemListStart
        \resumeItem{BUas Bachelor Scholarship recipient}
      \resumeItemListEnd
  \resumeSubHeadingListEnd

%-----------PROJECTS-----------
\section{Projects \& Technical Work}
  \resumeSubHeadingListStart

    \resumeProjectHeading
      {\textbf{CV Pipeline -- Plant Phenotyping Platform} $|$ \emph{5-person group project, BUas} $|$ \href{https://github.com/alex-krasnoshtanov/CV-Pipeline-Deployment-Platform}{\underline{Code}} $|$ \href{https://alex-krasnoshtanov.github.io/CV-Pipeline-Deployment-Platform/}{\underline{Docs}}}{Apr -- Jun 2026}
      \resumeItemListStart
        \resumeItem{Built an MLOps platform for plant organ segmentation and root-tip detection: U-Net (ResNet-34 encoder) served behind an authenticated FastAPI backend with a Next.js frontend, PostgreSQL and Prometheus}
        \resumeItem{Ran the full model lifecycle from Airflow DAGs -- data versioning and Azure ML training, drift detection on confidence, latency and error rate, and retraining driven by researcher feedback on flagged predictions}
        \resumeItem{Owned the deployment and CI/CD tier: one codebase shipping to Docker Compose, on-premise Portainer and Azure Container Apps, with blue-green cutover, health-gated rollback and OIDC-federated cloud auth; CI gates 85\% coverage across 430 tests}
      \resumeItemListEnd

    \resumeProjectHeading
      {\textbf{Dutch Sign Language Fingerspelling Trainer} $|$ \emph{Personal Project} $|$ \href{https://github.com/alex-krasnoshtanov/DSL-Learning}{\underline{Code}} $|$ \href{https://alex-krasnoshtanov.github.io/DSL-Learning/}{\underline{Docs}}}{May -- Sep 2026}
      \resumeItemListStart
        \resumeItem{Real-time NGT alphabet trainer: MediaPipe hand landmarks feed a PyTorch residual MLP for the 24 static letters (95\% validation accuracy) and a bidirectional LSTM over 30-frame sequences for the two letters requiring motion (92\%)}
        \resumeItem{Sole author: FastAPI WebSocket backend streaming to a Next.js 16 / TypeScript frontend, bilingual EN/NL, containerised with CI and published documentation}
      \resumeItemListEnd

    \resumeProjectHeading
      {\textbf{NPEC Plant Phenotyping System} $|$ \emph{Individual Project, BUas} $|$ \href{https://github.com/alex-krasnoshtanov/NPEC-Plant-Phenotyping}{\underline{Code}}}{Nov 2025 -- Jan 2026}
      \resumeItemListStart
        \resumeItem{Ranked \textbf{\#2 of 77} in Kaggle competition for plant root segmentation; selected to present results}
        \resumeItem{Built the full autonomous loop -- U-Net root segmentation, root-tip extraction, pixel-to-robot coordinate transform and OT-2 pipette control in a PyBullet simulation -- and benchmarked PID against a PPO reinforcement-learning controller}
      \resumeItemListEnd

    \resumeProjectHeading
      {\textbf{Emotion Analysis Pipeline} $|$ \emph{Group Project -- Dutch Media Client}}{Sep -- Nov 2025}
      \resumeItemListStart
        \resumeItem{Developed production NLP pipeline processing 40-minute video episodes end-to-end: Faster-Whisper V3 transcription with VAD, multilingual emotion classification (6 Ekman categories plus neutral), and automated visualisation}
        \resumeItem{Fine-tuned DeBERTa and Russian transformer models on merged datasets augmented with Mistral-generated synthetic data; cross-validated across languages to stabilise multilingual predictions}
      \resumeItemListEnd

    \resumeProjectHeading
      {\textbf{MARBET Event Support Chatbot} $|$ \emph{2-Person Sprint} $|$ \href{https://github.com/alex-krasnoshtanov/MARBET-Chatbot}{\underline{Code}}}{Apr -- May 2025}
      \resumeItemListStart
        \resumeItem{Built self-hosted RAG system (Llama 3.3 70B via Ollama, LangChain, FAISS) for corporate event Q\&A with multi-format ingestion and Gradio web interface; client deployed to production}
      \resumeItemListEnd

  \resumeSubHeadingListEnd

%-----------HACKATHONS & COMPETITIONS-----------
\section{Hackathons \& Competitions}
  \resumeSubHeadingListStart

    \resumeProjectHeading
      {\textbf{BrabantHack 2026} $|$ \emph{Winner -- Demcon Deep Tech Track} $|$ \href{https://github.com/alex-krasnoshtanov/Detection-by-Shadow}{\underline{Code}}}{Apr 2026}
      \resumeItemListStart
        \resumeItem{Won the Demcon track with a model that locates a pedestrian standing entirely outside the camera frame from the shadow they cast into it: \textbf{0.626 mean IoU} on the organisers' hidden test set, team of 3}
        \resumeItem{Drove the accuracy by replacing direct four-corner box regression with targets defined relative to the frame edge -- side, distance from edge, width, height, y-centre -- on a ResNet-50 trunk fused with 19 hand-crafted shadow descriptors; 0.468 to 0.610 mean IoU on a held-out split, reproducible from the repository}
      \resumeItemListEnd

    \resumeProjectHeading
      {\textbf{HERE Technologies Hackathon} $|$ \emph{RDM Verification Challenge}}{2026}
      \resumeItemListStart
        \resumeItem{Built XGBoost classification pipeline for Restricted Driving Manoeuvre verification combining probe GPS features with Mapillary and OSM geospatial signals; designed synthetic data scrambling for NDA-compliant delivery}
      \resumeItemListEnd

  \resumeSubHeadingListEnd

%-----------TECHNICAL SKILLS-----------
\section{Technical Skills}
 \begin{itemize}[leftmargin=0.15in, label={}]
    \small{\item{
     \textbf{Languages}{: Python, SQL, TypeScript} \
     \textbf{ML / Deep Learning}{: PyTorch, TensorFlow, Keras, Scikit-learn, XGBoost, HuggingFace Transformers} \
     \textbf{Computer Vision}{: OpenCV, U-Net, MediaPipe, PyBullet} \
     \textbf{NLP \& Audio}{: LangChain, Faster-Whisper, FAISS, NLTK, Ollama} \
     \textbf{MLOps \& Cloud}{: Azure ML, Azure Container Apps, Airflow, MLflow, Prometheus, GitHub Actions CI/CD} \
     \textbf{Infrastructure}{: Docker, Docker Compose, Portainer, FastAPI, Next.js, PostgreSQL, Git} \
     \textbf{Data \& Analytics}{: Pandas, NumPy, Power BI} \
     \textbf{Spoken Languages}{: English (fluent), Ukrainian (native), Russian (fluent)}
    }}
 \end{itemize}

\end{document}
% build trigger
